\documentclass[12pt,a4paper]{article}
\usepackage{ctex}          % 中文支持
\usepackage{amsmath}       % 数学环境
\usepackage{amssymb}       % 数学符号
\usepackage{geometry}      % 页面边距
\usepackage{booktabs}      % 三线表
\usepackage{array}         % 表格列格式
\usepackage{enumitem}      % 列表样式

\geometry{left=2.5cm,right=2.5cm,top=2.5cm,bottom=2.5cm}

\title{\heiti 拟合优度的度量——可决系数 \(R^2\) 的详细讲解}
\author{}
\date{}

\begin{document}

\maketitle

\section*{第一步：我们要解决什么问题？（引入概念）}

想象一下，我们在纸上画了很多散点（代表数据），然后我们画了一条直线试图穿过这些点。
这条直线就是我们的“样本回归线”，它代表了我们对数据关系的一种猜测或“拟合”。

\begin{itemize}[itemsep=0pt,topsep=0pt]
    \item \textbf{拟合}：就是画一条线来代表这些散点的趋势。
    \item \textbf{拟合优度}：这条线画得好不好？它和真实的点靠得有多近？这就是“拟合优度”要回答的问题。
\end{itemize}

为了用数字来衡量“好”的程度，我们需要一把尺子。这把尺子就是\textbf{可决系数（\(R^2\)）}。
而要理解这把尺子，我们必须先学会把数据的波动（变差）拆开来分析。

\section*{第二步：总变差的分解 —— 把数据的波动拆开来看}

我们先看一个核心公式（基于样本回归函数）：
\begin{equation}
    Y_i = \hat{Y}_i + e_i
\end{equation}
别被符号吓到，我们来解释每一个字母：
\begin{itemize}[itemsep=0pt,topsep=0pt]
    \item \(Y_i\)：\textbf{真实的观测值}。比如，第 \(i\) 个家庭真正花了多少钱。这是图上的那个\textbf{实心圆点}的高度。
    \item \(\hat{Y}_i\)：\textbf{估计值（预测值）}。这是我们画的回归线在第 \(i\) 个点位置上的高度，也就是线\textbf{预测}这个家庭应该花多少钱。
    \item \(e_i\)：\textbf{残差}。就是真实值 \(Y_i\) 和预测值 \(\hat{Y}_i\) 之间的距离差（\(e_i = Y_i - \hat{Y}_i\)）。如果线画得完美，这个差就是 0。
\end{itemize}

\subsection*{2.1 引入一个参考标准：平均值 \(\bar{Y}\)}

光看每个点离回归线多远还不够，我们需要一个比较的\textbf{基准}。这个基准就是所有真实值 \(Y_i\) 的\textbf{平均值 \(\bar{Y}\)}。

想象一个场景：你猜一个人的身高。
\begin{itemize}[itemsep=0pt,topsep=0pt]
    \item 如果没有任何信息（比如不知道体重），你猜的平均身高 \(\bar{Y}\) 就是最稳妥的猜测。
    \item 现在，你知道了体重 \(X\)，你用回归线算出了一个更精准的猜测 \(\hat{Y}_i\)。
    \item 我们要看的是：\textbf{用了体重信息（回归线）之后，你的猜测比瞎猜（用平均值）好了多少？}
\end{itemize}

\subsection*{2.2 关键的离差恒等式推导（极其详细，无跳步）}

我们要比较三个距离：
\begin{enumerate}[itemsep=0pt,topsep=0pt,label=(\arabic*)]
    \item \textbf{总离差}：真实值 \(Y_i\) 离开平均值 \(\bar{Y}\) 有多远？即 \((Y_i - \bar{Y})\)。
    \item \textbf{解释离差}：回归线预测值 \(\hat{Y}_i\) 离开平均值 \(\bar{Y}\) 有多远？即 \((\hat{Y}_i - \bar{Y})\)。
    \item \textbf{未解释离差}：真实值 \(Y_i\) 离开回归线预测值 \(\hat{Y}_i\) 有多远？即 \((Y_i - \hat{Y}_i) = e_i\)。
\end{enumerate}

\textbf{推导开始：}

我们已知回归线方程：\(Y_i = \hat{Y}_i + e_i\)。

\textbf{步骤 1}：等式两边同时减去平均值 \(\bar{Y}\)。
\[
Y_i - \bar{Y} = \hat{Y}_i + e_i - \bar{Y}
\]

\textbf{步骤 2}：把右边稍微重新排列一下，把 \(\hat{Y}_i\) 和 \(-\bar{Y}\) 放在一起。
\begin{equation}
    Y_i - \bar{Y} = (\hat{Y}_i - \bar{Y}) + e_i
\end{equation}

\textbf{结论}：这告诉我们，\textbf{总离差} 刚好等于 \textbf{回归线解释了的离差} 加上 \textbf{回归线没解释的残差}。
\begin{itemize}[itemsep=0pt,topsep=0pt]
    \item 如果回归线完全没用（\(\hat{Y}_i = \bar{Y}\)），那么解释离差为 0，总离差全是残差。
    \item 如果回归线完美（\(e_i = 0\)），那么总离差全被解释了。
\end{itemize}

\subsection*{2.3 为什么要把离差平方再加总？（\(\mathrm{TSS} = \mathrm{ESS} + \mathrm{RSS}\)）}

前面我们只看了\textbf{一个点}的情况。现在我们要看\textbf{所有点}的整体情况。

但是，直接加总离差 \((Y_i - \bar{Y})\) 会有一个致命问题：有的点比平均值高（正数），有的比平均值低（负数），一加就抵消了，总和接近 0，看不出波动大小。

解决办法：\textbf{平方！} 平方之后负数变正数，越远的点贡献越大。

我们对刚才的离差等式两边进行\textbf{平方，然后对所有的 \(i\)（从 1 到 \(n\) 个样本）求和}。

\textbf{等式：} \((Y_i - \bar{Y}) = (\hat{Y}_i - \bar{Y}) + e_i\)

\textbf{两边平方：}
\[
(Y_i - \bar{Y})^2 = [(\hat{Y}_i - \bar{Y}) + e_i]^2
\]

\textbf{展开右边（完全平方公式 \((a+b)^2 = a^2 + 2ab + b^2\)）：}
\[
(Y_i - \bar{Y})^2 = (\hat{Y}_i - \bar{Y})^2 + 2(\hat{Y}_i - \bar{Y})e_i + e_i^2
\]

\textbf{对所有样本求和（\(\sum_{i=1}^{n}\)）：}
\[
\sum (Y_i - \bar{Y})^2 = \sum (\hat{Y}_i - \bar{Y})^2 + 2\sum (\hat{Y}_i - \bar{Y})e_i + \sum e_i^2
\]

\textbf{关键跳跃点解释（这里绝对不停留）：} 在最小二乘法（OLS）的假设下，中间项 \(2\sum (\hat{Y}_i - \bar{Y})e_i\) \textbf{恰好等于 0}。
（因为残差和预测值的协方差为 0，这是 OLS 的数学性质，对于初学者只要记住：画回归线的方法保证了这一项消掉了）。

\textbf{最终简化结果：}
\begin{equation}
    \sum (Y_i - \bar{Y})^2 = \sum (\hat{Y}_i - \bar{Y})^2 + \sum e_i^2
\end{equation}

这就是 \textbf{式 (2.47)} 和 \textbf{(2.48)} 的来历。用简洁的代号表示就是：
\begin{equation}
    \mathrm{TSS} = \mathrm{ESS} + \mathrm{RSS}
\end{equation}

\begin{itemize}[itemsep=0pt,topsep=0pt]
    \item \textbf{TSS (Total Sum of Squares)}：总离差平方和。代表 \(Y\) 原本的波动有多大。
    \item \textbf{ESS (Explained Sum of Squares)}：回归平方和。代表\textbf{回归线能解释}的波动有多大。
    \item \textbf{RSS (Residual Sum of Squares)}：残差平方和。代表\textbf{回归线解释不了}的波动有多大。
\end{itemize}

\section*{第三步：可决系数 \(R^2\) —— 一把衡量好坏的尺子}

现在我们有了 TSS、ESS、RSS。我们想看回归线“解释了的波动”占总波动的\textbf{比例}。

\subsection*{3.1 公式推导}

既然 \( \mathrm{TSS} = \mathrm{ESS} + \mathrm{RSS} \)，我们两边同时除以 \(\mathrm{TSS}\)（假设 TSS 不为 0）：
\[
\frac{\mathrm{TSS}}{\mathrm{TSS}} = \frac{\mathrm{ESS}}{\mathrm{TSS}} + \frac{\mathrm{RSS}}{\mathrm{TSS}}
\]
\[
1 = \frac{\mathrm{ESS}}{\mathrm{TSS}} + \frac{\mathrm{RSS}}{\mathrm{TSS}}
\]

我们定义 \textbf{可决系数 \(R^2\)} 为 \(\frac{\mathrm{ESS}}{\mathrm{TSS}}\)。把它移项一下：
\begin{equation}
    R^2 = \frac{\mathrm{ESS}}{\mathrm{TSS}} = 1 - \frac{\mathrm{RSS}}{\mathrm{TSS}}
\end{equation}

\subsection*{3.2 直观理解 \(R^2\) 的大小意味着什么？}

\begin{itemize}[itemsep=6pt,topsep=0pt]
    \item \textbf{情况 1：拟合极好（\(R^2 = 0.99\)）}\\
    这意味着 \(R^2 = 0.99\)，也就是 \(\frac{\mathrm{ESS}}{\mathrm{TSS}} = 0.99\)。\\
    解释：总波动中，有 99\% 是因为 \(X\) 的变化引起的（被回归线抓住了），只有 1\% 是乱七八糟的噪声（残差）。\\
    \textbf{结论：} 线画得特别好，点几乎都贴在线上。

    \item \textbf{情况 2：拟合极差（\(R^2 = 0.01\)）}\\
    这意味着 \(\frac{\mathrm{RSS}}{\mathrm{TSS}} \approx 0.99\)。\\
    解释：总波动中，只有 1\% 被 \(X\) 解释，剩下的 99\% 都不知道哪来的（全在残差里）。\\
    \textbf{结论：} 这条线画了等于没画，还不如直接画一条平均值线。

    \item \textbf{情况 3：取值范围}\\
    因为 ESS 最大等于 TSS（完美拟合），最小等于 0（完全没用）。\\
    所以 \(R^2\) 永远在 \textbf{0 到 1} 之间。它不会出现负数（在带截距项的 OLS 回归中）。
\end{itemize}

\subsection*{3.3 书上的计算例子（一步步算）}

书里给了一个例子数据，虽然我们没看到表 2.4 的全貌，但给了两个关键结果：
\begin{itemize}[itemsep=0pt,topsep=0pt]
    \item 残差平方和 \(\mathrm{RSS} = \sum e_i^2 = 46116.02\)
    \item 总平方和 \(\mathrm{TSS} = \sum y_i^2 = 7561992\)
\end{itemize}

我们代入公式 \(R^2 = 1 - \frac{\mathrm{RSS}}{\mathrm{TSS}}\)：

\textbf{步骤 1：} 算除法 \(\frac{46116.02}{7561992} \approx 0.006098\)（约等于千分之六）。

\textbf{步骤 2：} 用 1 减去它 \(1 - 0.006098 = 0.993902\)。

\textbf{步骤 3：} 化为百分数 \(99.39\%\)。

这就像一个考试，满分 100 分，你只扣了 0.6 分，考了 99.4 分。非常棒！

\section*{第四步：可决系数 \(R^2\) 与相关系数 \(r\) 的关系（极其详细的代数推导）}

这是最容易混淆的地方。书上说“在一元线性回归中，\(R^2 = r^2\)”。我们不仅要记住这句话，还要看懂是怎么变出来的。

\subsection*{4.1 先回忆相关系数 \(r\) 的公式}

\begin{equation}
    r = \frac{\sum (X_i - \bar{X})(Y_i - \bar{Y})}{\sqrt{\sum (X_i - \bar{X})^2 \sum (Y_i - \bar{Y})^2}}
\end{equation}
（为了书写方便，我们记分子为 \(\sum xy\)，分母为 \(\sqrt{\sum x^2 \sum y^2}\)）

\subsection*{4.2 推导 \(R^2\) 的另一副面孔}

我们已经知道 \(R^2 = \frac{\mathrm{ESS}}{\mathrm{TSS}} = \frac{\sum (\hat{Y}_i - \bar{Y})^2}{\sum (Y_i - \bar{Y})^2}\)。

关键的一步是处理分子 \(\mathrm{ESS}\)。在一元回归中，我们知道回归线的斜率 \(\hat{\beta}_2\) 的公式是：
\begin{equation}
    \hat{\beta}_2 = \frac{\sum (X_i - \bar{X})(Y_i - \bar{Y})}{\sum (X_i - \bar{X})^2}
\end{equation}

另外，根据回归线的性质，预测值 \(\hat{Y}_i = \hat{\beta}_1 + \hat{\beta}_2 X_i\)，且有 \(\hat{Y}_i - \bar{Y} = \hat{\beta}_2 (X_i - \bar{X})\)。\\
\textbf{（暂停理解一下：} 回归线经过中心点 \((\bar{X}, \bar{Y})\)，所以预测值偏离平均值的幅度，等于斜率 \(\hat{\beta}_2\) 乘以 \(X\) 偏离其平均值的幅度。\textbf{）}

现在我们来化简 \(\mathrm{ESS}\)（分子）：
\[
\mathrm{ESS} = \sum (\hat{Y}_i - \bar{Y})^2
\]
把 \(\hat{Y}_i - \bar{Y} = \hat{\beta}_2 (X_i - \bar{X})\) 代入：
\[
\mathrm{ESS} = \sum [\hat{\beta}_2 (X_i - \bar{X})]^2
\]
\[
\mathrm{ESS} = \sum \hat{\beta}_2^2 (X_i - \bar{X})^2
\]
常数 \(\hat{\beta}_2^2\) 可以提到求和符号外面：
\[
\mathrm{ESS} = \hat{\beta}_2^2 \sum (X_i - \bar{X})^2
\]

\textbf{现在把 \(\hat{\beta}_2\) 的公式代进去：}
\[
\mathrm{ESS} = \left( \frac{\sum (X_i - \bar{X})(Y_i - \bar{Y})}{\sum (X_i - \bar{X})^2} \right)^2 \times \sum (X_i - \bar{X})^2
\]

注意看括号外面的指数。外面是平方，里面分母有 \(\sum (X_i - \bar{X})^2\)，分子有 \(\sum (X_i - \bar{X})^2\) 的一次方。约分掉一个：
\begin{equation}
    \mathrm{ESS} = \frac{[\sum (X_i - \bar{X})(Y_i - \bar{Y})]^2}{\sum (X_i - \bar{X})^2}
\end{equation}

\subsection*{4.3 最后一步：凑出 \(R^2\) 的样子}

我们已经把分子 ESS 化简好了。现在看整个 \(R^2\)：
\[
R^2 = \frac{\mathrm{ESS}}{\mathrm{TSS}} = \frac{\frac{[\sum (X_i - \bar{X})(Y_i - \bar{Y})]^2}{\sum (X_i - \bar{X})^2}}{\sum (Y_i - \bar{Y})^2}
\]

这等于：
\begin{equation}
    R^2 = \frac{[\sum (X_i - \bar{X})(Y_i - \bar{Y})]^2}{\sum (X_i - \bar{X})^2 \cdot \sum (Y_i - \bar{Y})^2}
\end{equation}

看看这个式子，它正好就是 \textbf{相关系数 \(r\) 的公式} 的 \textbf{平方}！
因为：
\begin{equation}
    r^2 = \left( \frac{\sum (X_i - \bar{X})(Y_i - \bar{Y})}{\sqrt{\sum (X_i - \bar{X})^2 \sum (Y_i - \bar{Y})^2}} \right)^2 = \frac{[\sum (X_i - \bar{X})(Y_i - \bar{Y})]^2}{\sum (X_i - \bar{X})^2 \sum (Y_i - \bar{Y})^2}
\end{equation}

\textbf{结论：} 在一元回归里，\(R^2 = r^2\)。

\subsection*{4.4 为什么它们概念不同？（这部分极其重要，防止用错）}

虽然数值上平方一下就是对方，但它们的\textbf{灵魂}是完全不同的：

\begin{table}[h]
    \centering
    \renewcommand{\arraystretch}{1.3}
    \begin{tabular}{p{2.2cm}<{\centering} p{5.5cm} p{5.5cm}}
        \toprule
        \textbf{对比项} & \textbf{可决系数 \(R^2\)} & \textbf{相关系数 \(r\)} \\
        \midrule
        方向性 & 有方向的因果关系。比如 \(X\) 是原因，\(Y\) 是结果。\(R^2\) 衡量的是 \(X\) 解释了 \(Y\) 变动的百分之几。 & 无方向的相互依存关系。\(X\) 和 \(Y\) 一起变，变动的紧密程度如何。 \\
        \midrule
        正负号 & 永远非负 \(0 \leq R^2 \leq 1\)。因为它衡量的是比例（面积/面积）。 & 可正可负 \(-1 \leq r \leq 1\)。正号表示同向变化，负号表示反向变化。 \\
        \midrule
        比喻 & 用身高（\(X\)）去预测体重（\(Y\)），身高能解释体重变动的 70\%。你\textbf{不会}反过来说体重解释身高变动的 70\%。 & 身高和体重的关联度是 0.8。这是一个对称的数值，不分谁解释谁。 \\
        \bottomrule
    \end{tabular}
\end{table}

\section*{总结}

整篇文章的逻辑链条如下：
\begin{enumerate}[itemsep=0pt,topsep=0pt]
    \item \textbf{拆开波动}：把每个点的差距拆成“回归线解释的”和“没解释的”。
    \item \textbf{平方求和}：为了避免正负抵消，平方后求和得到 \(\mathrm{TSS} = \mathrm{ESS} + \mathrm{RSS}\)。
    \item \textbf{计算比例}：看 \(\mathrm{ESS}\) 占 \(\mathrm{TSS}\) 的多少，这就是 \textbf{\(R^2\)}。
    \item \textbf{对比概念}：在一元情况下，\(R^2\) 正好等于相关系数 \(r\) 的平方，但理解上一个是因果解释程度，一个是相互关联程度。
\end{enumerate}

希望这样的讲解能够让你完全明白每一步是在做什么。

\end{document}